\section{Discussion}\label{sec:discussion}
The measurements of Section~\ref{sec:evaluation} and the security analysis of Section~\ref{sec:security} together answer the question this thesis set out to ask. That question is whether a zkVM is a viable substrate for post-quantum anonymous credentials, and at what cost, not how fast the zkVM runs in isolation. We begin with the main finding, the anonymity--deployability tension, then give the cost results as guidance for a deployer.

\subsection{The Anonymity--Deployability Tension}\label{ssec:tension-discussion}
The anonymity--deployability tension of Section~\ref{sec:security} matters more than any single timing figure. Two things hold together on the target hardware at $\lambda=80$. First, the proof that provides anonymity both runs and binds. Both credential relations commit their public statement to the receipt, the full BDEC lifecycle runs end to end (two cross-issuer issuances and a $k=2$ showing), and the zero-knowledge wrap that BDEC's anonymity requires completes within the $24$~GB budget. The standalone verification wraps in $46.83$~minutes and the \emph{statement-bound} issuance receipt in $4.18$~hours, each near $15$~GiB (Section~\ref{sec:evaluation}). So binding and zero-knowledge combine in one measured proof, not two separate demonstrations. Second, that wrap is pairing-based, so the anonymity it delivers is classical only. A deployer who needs \emph{post-quantum} anonymity on this stack still cannot get it, and the reason is not a shortage of memory or time. A reduced verifying-key table built from the guest's own shard shapes (\S\ref{ssec:resource-frontier}) closes the chain without a multi-day rebuild, so the memory cost of that rebuild is not what blocks anonymity. What blocks it is the primitive. The only wrap the toolchain provides uses a pairing again. This is the single remaining side of the dual obstruction, the second obstruction (Section~\ref{sec:security}), a limit of primitives rather than of resources. Recovering post-quantum anonymity needs a post-quantum-sound zero-knowledge wrap, a STARK- or lattice-based recursive argument the toolchain does not yet provide.

The obstruction is specific to the zkVM stack. The static-circuit regime does produce a zero-knowledge post-quantum proof of the Loquat single-verification building block in about $22$~minutes (mean of three runs, $21.4$--$22.9$) on the same machine (\rr{docs/measurements/audit5_campaign_20260612}), at the per-change regeneration cost weighed below. And the classical anonymity the zkVM stack does deliver is itself conditional. It rests on two open assumptions of Section~\ref{sec:security}: PLUM's signature-transcript simulatability (Assumption~\ref{asm:szk}) and, for issuance, its signature key-privacy (Assumption~\ref{asm:keypriv}), now essential because statement-binding makes the issuance credential a public input, together with the zero-knowledge wrap that BDEC's anonymity requires. The tension holds independently of how favourable the timing numbers are and of the open AIR-soundness Assumption~\ref{asm:air}, which conditions the feasibility figure rather than the obstruction.

\subsection{A Substrate-Selection Rule}\label{ssec:decision}
The two regimes do not dominate each other. Each trades a different cost. The right choice depends on one property of the deployment: how often the verification rule changes compared with how often proofs are generated.

\begin{itemize}
  \item \textbf{When the predicate is fixed at deployment time} and proofs are generated frequently, the static-circuit regime is preferable: its specialised R1CS amortises across many showings, and it pays its update cost only once. The zkVM's per-proof overhead (Section~\ref{sec:evaluation}) is not justified when there is nothing to update.
  \item \textbf{When the predicate is chosen at presentation time, or the signature or security parameters evolve over the deployment's lifetime}, the static-circuit regime pays its update-churn cost (Table~\ref{tab:churn}) \emph{repeatedly}, and the zkVM regime, which absorbs each such change as a guest-program edit, becomes the structurally appropriate substrate, at the proving cost quantified in Section~\ref{sec:evaluation}. The exception is a change to the precompiled primitive itself, which regenerates its chip (below). For a non-preprocessing static argument such as Aurora this update-churn cost is a regeneration-and-re-audit (governance and assurance) burden rather than a large wall-clock penalty, which the next paragraph makes precise.
\end{itemize}

\noindent This rule can be made precise. Over a deployment lifetime of $P$ proof generations and $U$ localised predicate changes, the amortised cost of each regime is, schematically,
\begin{align}
  \text{static:}\quad & U\cdot R_{\mathsf{static}} + P\cdot t^{P}_{\mathsf{static}},
  \label{eq:static-cost}\\
  \text{zkVM:}\quad & U\cdot 0 + P\cdot t^{P}_{\mathsf{zkVM}},
  \label{eq:zkvm-cost}
\end{align}
where $R_{\mathsf{static}}$ is the per-change recompilation cost and $t^{P}$ the per-proof time; the zkVM absorbs each predicate change as a guest-program edit, so its update term is zero. The static regime is cheaper only while the predicate-change rate stays below the crossover
\begin{equation}
  \Bigl(\tfrac{U}{P}\Bigr)^{\!*} \;=\; \frac{t^{P}_{\mathsf{zkVM}}-t^{P}_{\mathsf{static}}}{R_{\mathsf{static}}},
  \label{eq:crossover}
\end{equation}
above which the zkVM is the cheaper substrate: the threshold is the ratio of the zkVM's per-proof penalty to the static circuit's per-change recompilation cost. The recompilation cost is measured (Section~\ref{sec:evaluation}): for a transparent non-preprocessing argument such as Aurora it is $R_{\mathsf{static}}\approx 0.3$~s (R1CS construction plus a microsecond parameter setup, single timed runs), and it sits outside the multi-minute prover, which is paid per proof. At that magnitude, and with the zkVM dearer per proof on the cross-scheme reading available to us, the wall-clock crossover would favour the static circuit in any realistic churn regime. Against Aurora, then, the zkVM has no wall-clock flexibility advantage on that reading. The same-scheme sign of the crossover is settled below as a conjecture, not a measurement. The flexibility claim we can defend is instead the \emph{qualitative} regeneration-and-re-audit footprint of Table~\ref{tab:churn} (zero on the zkVM side, $\geq 1$ on the static side per change), an engineering-and-assurance cost paid by humans, not prover time.

\noindent A wall-clock advantage appears only against a \emph{preprocessing} baseline, which Section~\ref{sec:evaluation} measures for the transparent preprocessing argument Fractal (Table~\ref{tab:rstatic-fractal}) and which would hold in principle for Groth16~\cite{groth2016}, Marlin~\cite{chiesa2020marlin}, or circuit-specific PLONK~\cite{gabizon2019plonk}; the non-preprocessing Aurora that BDEC uses does not offer one.

\noindent This zero-regeneration property is partitioned rather than absolute. A change confined to the verification predicate is a guest-program edit that reuses the existing universal relation, but a change to the hash family or to the field width still regenerates and re-audits the Griffin-$\mathbb{F}_p^{192}$ precompile, so the decision rule applies to predicate churn and not to substrate-level changes of the signature itself.

\noindent What the cost model licenses is limited: it fixes the \emph{form} of the rule and the \emph{sign} of every term, but not the numeric location of the crossover for PLUM. So the rule gives directional guidance, telling a deployer which side a high- or low-churn deployment lies on without yet handing them a numeric threshold to evaluate.

\noindent One asymmetry the rule abstracts away: the zkVM's per-proof cost is flat only on the issuance side. A $k$-credential presentation's cost grows with the credential count: execute-mode cycles rise from $1.36\times10^{11}$ at $k=1$ to $1.82\times10^{11}$ at $k=2$, and the zero-knowledge wrap's core-shard count and peak memory grow with them (both the $k=1$ and $k=2$ wraps measured, within the $24$~GB budget; Section~\ref{sec:evaluation}), so for a large number of shown credentials the zkVM showing cost projects past a practical bound even where the churn rate favours the zkVM. The credential count is therefore a second axis the deployer weighs alongside the churn rate the rule captures.

\noindent The contribution of this thesis to that rule is the \emph{quantification of both sides}: the zkVM's per-proof cost is measured (Tables~\ref{tab:plum-verify}--\ref{tab:bdec}), and the structural form of the static circuit's update cost is enumerated (Table~\ref{tab:churn}). The static circuit's recompile \emph{time} is measured ($R_{\mathsf{static}}\approx 0.3$~s, \S\ref{ssec:churn-eval}) and is sub-second because Aurora is non-preprocessing; the remaining term is a \emph{same-scheme} per-proof prover time on both substrates, and this we do not measure: neither Loquat nor PLUM is run on both substrates at matched security and rate. We expect the specialised static circuit to be the cheaper per-proof prover, since a native constraint system avoids the zkVM's emulation of the signature's field arithmetic, but the sign of that crossover is a conjecture rather than a measurement. Settling it requires an $\mathbb{F}_p$ ($199$-bit) Aurora harness to run PLUM statically (the current harness is $127$-bit only), which we identify as the required follow-up experiment.

\subsection{What Generalises and What Does Not}\label{ssec:generalise}
The field-mismatch tax (Section~\ref{sec:theoretical}) and the update-churn dimension (Definition~\ref{def:update-churn}) are substrate-independent: they apply to any large-field algebraic-hash post-quantum signature of this form verified in any small-field zkVM that exposes application-defined precompiles, not only to PLUM in SP1. The specific magnitudes (Section~\ref{sec:evaluation}) do not generalise: they are tied to PLUM's $199$-bit field, the M5~Pro's $24$~GB envelope, and the CPU-bound prover of the SP1 zkVM studied. We drew the boundary deliberately, so that the shape of the cost carries to other pairings of a large-field signature with a small-field zkVM while the numbers stay pinned to the configuration measured here.

\subsection{Comparison to Prior Evaluations}\label{ssec:prior-eval}
As anticipated in Section~\ref{sec:related}, the prior-evaluation record omits two axes: the cost of a \emph{change} to the verification predicate, and the consumer-hardware proving frontier for an algebraic-hash post-quantum signature. We do not restate that gap here; instead we report what our measurements add to the record. For predicate-change cost, Table~\ref{tab:churn} fixes its structural form (zero regenerations on the zkVM side against $\geq 1$ per change on the static side) and \S\ref{ssec:churn-eval} bounds its wall-clock magnitude; for the frontier, \S\ref{ssec:resource-frontier} locates the binding constraint at the zero-knowledge wrap rather than the base proof: as built and measured the wrap completes within budget, so what binds is its post-quantum soundness and not its feasibility (the base-proof cost is the honest-path cost of the chip and carries Assumption~\ref{asm:air}). Where our per-proof timings overlap the existing record, they agree in order of magnitude with the hitherto-qualitative observation that algebraic-hash workloads are expensive inside small-field zkVMs, now quantified at the verification level (Tables~\ref{tab:plum-verify}--\ref{tab:ablation}). The next section combines these into the thesis's overall conclusion.
